\section{Corpus and Method}
\label{sec:methodology}

\subsection{The initial plan, and what changed}

The plan at the start was to look for complete cross-agent escalation chains
in public infrastructure as code on two platforms, \agentcore and \foundry.
Google Cloud was left out from the beginning. Two parts of that plan turned
out differently.

The first was the expectation that complete chains would be there to count.
None were. The question became how far the incomplete paths fall short, and
which single permission holds them back.

The second was the expectation that \foundry agents would be written into
templates the way other Azure resources are. They are not. Creating an agent
is a data plane call, and the identity that runs a template does not hold it.
The \foundry part became a measurement of how connections are set to
authenticate, together with a live lab, rather than a count of declared
agents.

The collection script for the \foundry corpus fetched Terraform files but not
Bicep files, and most \foundry material is written in Bicep. The corpus was
collected again with the file pattern widened.

\subsection{What we collected}

We collect public repositories that deploy the platforms in question, in three
families: Terraform projects that declare \agentcore agent runtimes,
CloudFormation and CDK deployments of the same service, and \foundry
templates.

Every collection query records whether it hit the search API's result ceiling.
No query in either corpus did, so for the queries we issued the collection is
exhaustive rather than a sample of what was available. A deployment using none
of our query terms is not in the corpus.

\subsection{What we excluded, and why}

A file under a test fixture, a snapshot directory, or a vendored dependency is
not a deployment, and a template byte identical to another is counted once.
The rule applies everywhere a count appears, including the role sharing
census. Every excluded unit is named in the artifacts.

\subsection{Three denominators}
\label{sec:populations}

Three different denominators appear in this paper and they are not the same
set.

\begin{description}
\item[Analysed projects (\nProjects).] Terraform projects that declare at
  least one agent runtime and that the parser read without failing. This is
  the denominator for escalation results.
\item[Chain denominator (\nChainDenominator).] The deployable distinct units
  the chain analysis actually covered: \nChainTerraform Terraform projects and
  \nChainCfn CloudFormation templates. \nChainNotCovered units in the
  population are not covered, and they are named rather than assumed clean.
\item[Foundry directories (\nFoundryDirs).] Template directories in the
  \foundry collection. Only \nFoundryAuthDirs of them declare a connection
  resource, and that smaller set is the denominator for the \api{authType}
  figures.
\end{description}

Figure~\ref{fig:pipeline} shows how the collection narrows to each of them.

\begin{figure}[t]
\centering
\begin{tikzpicture}[
  font=\scriptsize,
  box/.style={draw, rounded corners=2pt, align=center, inner sep=3pt,
              text width=54mm, minimum height=6mm, line width=0.5pt},
  keep/.style={box, draw=cIdentity, fill=cIdentity!8},
  narrow/.style={box, draw=cNeutral, fill=cFaint},
  final/.style={box, draw=cAccent, fill=cAccent!8},
  down/.style={-{Stealth[length=1.8mm]}, line width=0.6pt},
  side/.style={font=\tiny, align=left, text width=20mm, inner sep=1pt},
]

\node[keep] (a) {\textbf{\nReposCollected} public repositories collected, no
  query hit the API ceiling};
\node[narrow, below=4mm of a] (b) {\textbf{\nReposWithRuntime} of them declare
  an agent runtime, holding \textbf{\nTemplateDirs} template directories};
\node[narrow, below=4mm of b] (c) {after removing \nTemplateExcluded test
  fixtures, snapshots and vendored copies, and \nTemplateForkCopies fork
  copies};
\node[final, below=4mm of c] (d) {\textbf{\nProjects} analysed Terraform
  projects\\\scriptsize denominator for escalation results};
\node[final, below=13mm of d] (e) {\textbf{\nChainDenominator} units covered by
  the chain analysis\\\scriptsize \nChainTerraform Terraform $+$ \nChainCfn
  CloudFormation};

\draw[down] (a) -- (b);
\draw[down] (b) -- (c);
\draw[down] (c) -- (d);
\draw[down] (d) -- node[side, right=1mm]
  {\nChainNotCovered units not covered, named individually} (e);

\end{tikzpicture}
\caption{How the collection narrows to each denominator. The two shaded boxes
are different populations and the paper does not mix them: a rate over
\nProjects and a rate over \nChainDenominator answer different questions.}
\label{fig:pipeline}
\end{figure}


\subsection{Verification}

Numbers come from artifacts, the artifacts are rebuilt by one script, and a
verification pass re-derives every published figure from them.
Appendix~\ref{sec:repro} gives the commands. Where a second pass re-derived a
number from the raw source, Appendix~\ref{sec:claims} reports its precision;
where none did, it says so.
